\documentclass[11pt]{article}
\usepackage[utf8]{inputenc}
\usepackage[T1]{fontenc}
\usepackage{geometry}
\usepackage{enumitem}
\usepackage{xcolor}
\usepackage{hyperref}
\usepackage{amsmath}

\geometry{margin=1in}

\newcommand{\resolved}[1]{\textcolor{teal}{\textbf{[RESOLVED]} #1}}
\newcommand{\open}{\textcolor{red}{\textbf{[OPEN]}} }
\newcommand{\partialr}{\textcolor{orange}{\textbf{[PARTIAL]}} }
\newcommand{\cref}[1]{\textcolor{purple}{\textbf{[#1]}} }

\title{\textbf{Singular Learning Theory for Applied Deep Learning}\\[8pt]
\large Needs Document: Open Items by Section}
\author{Working Document, aligned with goals (v5) and analysis (v3)}
\date{\today}

\begin{document}
\maketitle

\noindent\emph{Reconstructed from the flags in the analysis plan plus
items opened or resolved in the framing rounds. Every [NEED] and
[VERIFY] in the analysis plan points here. Status: \open{} not started;
\partialr{} evidence gathered, statement not pinned;
\textcolor{teal}{\textbf{[RESOLVED]}} citable, with source. Items
backing a load-bearing contingency in the goals document carry its tag
(\cref{C1}--\cref{C7}); resolving a tagged item settles the
corresponding contingency, and where the goals ledger records an upside
action, it is repeated here so the verification work is scoped knowing
what it purchases.}

% ============================================================================
\section{Section 1 Needs}
% ============================================================================

\subsection*{Pillar 1 --- wrong results}
\begin{itemize}[itemsep=4pt]
\item \open \cref{C2} Wei--Lau singular VI (arXiv:2302.06035): pin the
precise statement of mean-field's failure on singular posteriors and of
the resolution-of-singularities family's fix; identify the reproducible
small experiment for the \S1 figure; \emph{and} check the scale at which
singular VI is computable relative to the reader's models --- the scale
check decides \S4.3's register (provocative if computable at relevant
scale; humble geometry-settled/computation-frontier otherwise, per
framing decision (i)). If the statement check fails, Pillar 1 leans on
the degenerate-Hessian Laplace exhibit as anchor.
\item \open Normalization / scale-invariance: the van Laarhoven lineage;
the precise sense in which the weight-decay complexity rationale is void
under $f(x,\alpha\theta)=f(x,\theta)$; exact citations.
\end{itemize}

\subsection*{Pillar 2 --- misaimed fixes (and the audit table)}
\begin{itemize}[itemsep=4pt]
\item \partialr Koh--Liang influence functions: assumption (strict
convexity, PD Hessian) and the damping repair ($H+\lambda I$, $\lambda >
|E_0|$) are confirmed, including the observation that optimization leads
to critical points with a majority of near-zero eigenvalues. Remaining:
pin the primary locus of the ``damping large enough to distort'' result
(Basu et al.\ 2020 and the OOD-IF literature) with page/claim.
\item \partialr Second-order damping ritual: Martens 2010 damping trick
confirmed in secondary sources; cite from the primary. K-FAC damping:
pull the exact treatment from Martens--Grosse 2015.
\item \open Sharpness lineage: exact Dinh et al.\ 2017 statement; one
citable instance of the field repairing the sharpness measure
post-critique (e.g.\ adaptive/invariant sharpness variants) for the
``repaired the measure, not the object'' claim.
\item \open Audit table reconciliation (updated to goals v5): update the
drafted table (dl\_section1\_audit\_table.tex) to (a) \emph{remove} the
convergence-guarantees row from the table --- it is a level-2 item and
moves to \S4.2's recast (goals decision (k)); (b) add the third gate:
every row makes a parameter-space claim (posterior, uncertainty region,
influence, curvature, parameter count); (c) the masking framing in the
caption; (d) a diagnosis-tag column (geometry / sample size /
computation) per row.
\end{itemize}

\subsection*{Pillar 3 --- regularizing instead of redesigning}
\begin{itemize}[itemsep=4pt]
\item \open \cref{C5} Confirm the one-reflex framing holds across the
table's ``shipped repair'' column --- that damping, Gaussian families,
and sharpness penalties are genuinely one move at different sites and
can be presented as such without overreach. Now gated to
parameter-space rows only (easier, per the table's third gate).
Fallback if it fails: ``a family of regularity-restoring moves''; the
escalation ordering survives.
\end{itemize}

\subsection*{Pillar 4 --- folklore in the gap}
\begin{itemize}[itemsep=4pt]
\item \partialr Convergence guarantees, scoped (replaces the old
``guarantees assume non-degenerate Hessian'' item, which is \emph{not}
accurate as stated): modern nonconvex Adam/RMSProp analyses deliver
convergence to $\epsilon$-stationary points under (generalized)
smoothness and do not require a PD Hessian; guarantees promising minima
invoke PL or strong convexity; PL tolerates non-isolated minima but is
blind to strata. Remaining: select the two or three canonical citations
(e.g.\ relaxed-assumption Adam analyses; a PL-condition reference) and
one instance of tuning lore bridging the gap (warmup/schedule folklore).
\emph{Placement update (goals v5, decision (k)): the convergence-lore
folklore face moves from \S1 to \S4.2; \S1's folklore faces are the
variational-gap misattribution and the damping ritual only. The
citations gathered here serve \S4.2.}
\item \open Hagiwara: AIC fails to size perceptrons, traced to singular
structure --- confirm citation and year (referenced via Amari-school
secondary sources; find the primary).
\item \open The damping-ritual folklore face: assemble the
hand-tuned-$\lambda$ evidence across Laplace / influence / natural
gradient / K-FAC (rules of thumb, per-problem tuning) --- two or three
citable instances suffice.
\item \resolved{Watanabe capstone (the language had to be imported):
Watanabe's own exposition states that in deep learning a statistical
sub-model corresponds to an algebraic variety, that singularities are
made normal crossing via Hironaka's resolution, and that learning is
captured by two birational invariants (RLCT, singular fluctuation).
Source: Watanabe's SLT exposition and the 2009 Cambridge book. Supports
\S1.4 and \S2.2.}
\item \partialr Information-geometry precedent (the ``seen'' half of
seen-but-not-mainstreamed): Fukumizu 1996 (regularity condition of the
information matrix); Fukumizu--Amari 2000 (plateaus from singular
structure); Amari--Park--Ozeki 2002/2006 (neuromanifold singularities;
Milnor attractors); singular structure identified as common to MLPs, RBF
nets, mixtures, ARMA. Gathered from the literature pass; pin exact
citations in the bibliography. Remaining: keep the Watanabe--Amari
personal lineage OUT unless a primary source is found --- the claim stays
institutional.
\end{itemize}

% ============================================================================
\section{Section 2 Needs}
% ============================================================================

\begin{itemize}[itemsep=4pt]
\item \open Catalogue of architectural continuous symmetries with
citations (normalization scale-invariance; homogeneous-activation
rescaling; softmax shift; attention-head permutation as the discrete
part).
\item \open DGP degeneracy: state as presumption, not theorem
(``delete a node and the fit degrades'' is probably false); find the
honest citable form.
\item \resolved{Watanabe uses ``normal crossing'' and ``algebraic
variety in the parameter space'' for DL sub-models --- confirmed in his
own exposition; supports the \S2.2 thesis and the paper's positioning.}
\item \open The microscope FFN: fix the architecture (small enough for
2D KL slices; rich enough for a genuine crossing); produce the
hyperbolic-vs-parabolic contour figure and the contour
$\leftrightarrow$ eigenspectrum calibration figure.
\item \open \S2.4 promoted experiment (load-bearing for the sample-size
arm of the diagnosis): posterior/sample cloud at multiple $n$;
null-space profile vs.\ ($n \times$ degeneracy knob). The figure must
support both readings: small $n$ masks the geometry (regular tools
``work''); growing $n$ reveals it (conclusions may not transfer). The
figure's specified job now includes the diffusion mechanism (goals v5,
\S2 overview): the level sets fix how much room finite-$n$ estimates
have to wander; growing $n$ closes the room, revealing the set --- show
the cloud shrinking onto the crossing.
\item \open \cref{C4} \S2.4 literature check (higher bar now): the
$n$-revelation and $n$--dimension compounding claim against
finite-sample and dimension-dependence results in SLT before any strong
form is stated --- which also asks whether the compounding is genuinely
unspelled-out (a contribution to name) or already stated (cite and
lean). Fallback: shown-not-interpreted in \S2; strong form posed as
open in \S5. \emph{Upside if verified: restore the compounding to the
goals document's NOT bullet 1 contribution list and to walk-away item
5's lens sentence.}
\item \open Optional figure candidate (from the framing round): level-set
curvature concentration as an empirical crossing signature --- the
second fundamental form of the KL level sets is flat along single strata
and concentrates near stratum intersections; plottable on the microscope
FFN. A candidate, not a commitment; if used, it visualizes ``hyperbolic
level sets as the empirical signature'' without introducing curvature
language into the prose (the invariant content stays volume-based).
\end{itemize}

% ============================================================================
\section{Section 3 Needs}
% ============================================================================

\begin{itemize}[itemsep=4pt]
\item \open Watanabe's exact singular-model definition and analyticity
conditions (model vs.\ $K$); locations in Watanabe 2009/2013.
\item \open \cref{C6} Transformer activations are analytic: confirm
GELU/SiLU analyticity and their dominance in the workbench model class;
this carries decision (t)'s ``narrowing the gap for the reader's own
models'' sentence. Fallback: the sentence is cut; the caveat stands on
the microscope's rigor and the softplus-limit framing alone.
\item \open Explicit blow-up for the microscope FFN crossing;
before/after visualization (tangled hyperbolic $\to$ rectangular level
sets).
\item \open Hironaka statement with the conditions actually used.
\item \open Main theorem statement with conditions and limits; the
free-energy expansion including the $(m-1)\log\log n$ term; the tempered
posterior / WBIC $\beta = 1/\log n$ statement.
\item \open Compute $\lambda$ for the microscope FFN (needed for \S2.4's
concentration rate and \S3.5's validation).
\item \open \S3.5 validation run: theoretical free energy vs.\
thermodynamic integration vs.\ parameter-counting baseline at multiple
$n$; report the $n$ needed honestly (doubles as a computation-maturity
datapoint for \S4).
\end{itemize}

% ============================================================================
\section{Section 4 Needs}
% ============================================================================

\begin{itemize}[itemsep=4pt]
\item \open \textbf{The tools-and-maturity audit (gating task for \S4).}
Per tool: problem solved / diagnosis tag (geometry, sample size,
computation) / concept-maturity / computation-maturity / scale /
progress / rough edges / reader layer. Roster at minimum: WBIC-style
criteria, LLC/rLLC estimation, singular VI, susceptibilities, BIF,
sampler configurations in \texttt{devinterp}. Scale is an explicit
column: it assigns the register per decision (i).
\item \open \textbf{The practice inventory (second gating task, goals
decision (o)).} Per practitioner activity: claim made / tool used /
where the computation lives / the question it answers vs.\ the question
asked / verdict. Feeds the tool-by-question scope table.
\item \open Watanabe 2010: re-confirm the exact CV$\approx$WAIC statement
and gen$+$CV$\to 2\lambda/n$; the precise AIC/DIC failure statement.
\item \open \cref{C3} The small-$n$ selection statement (the
quantitative half of the two-questions hinge). Verified qualitatively:
CV/WAIC unbiased for generalization but inconsistent for selection; free
energy consistent, not unbiased; Watanabe's own deferral of the
selection comparison. Open: pin the quantitative small-$n$ claim
(selection signal shrinking against CV noise) in Watanabe's 2021 mixture
comparison or Yamazaki--Kaji 2013, or demonstrate it with the mixture
experiment (C1a below). Fallback: \S4.1 states the consistency asymmetry
without the finite-$n$ rate; the mixture experiment is reported as
illustration, not theorem.
\item \open \cref{C1} The conclusion-flip exhibit (backs the masking
thesis's ``structurally wrong where checkable''). One validated case
where a regular tool and its singular counterpart disagree on a
conclusion and the singular answer is verifiably right. Candidates,
cheapest first: (a) the small-$n$ mixture selection experiment (CV
selects the larger model; free energy selects the true one) --- doubles
as the cross-paper bridge to the applied-stats binomial mixture; (b) the
LLC optimizer exhibit if matched fit holds (C7); (c) BIF-vs-damped-IF
reordering of influential examples, validated by retraining. Fallback:
the alarm rests on the structural-wrongness exhibits; the
conclusion-flip moves to \S5 as the open empirical question; aim (i)'s
``structurally wrong where checkable'' softens to ``can be structurally
wrong.''
\item \open \cref{C7} Matched fit in the LLC optimizer exhibit (Lau et
al.): the $\lambda$ separation between SGD and entropy-SGD endpoints and
its low variance are verified as stated by the authors; \emph{not yet
verified}: that the endpoints have comparable training and test
performance (1.9M-parameter FFN on MNIST; 80 repetitions). The
matched-fit half carries the emblematic framing (``equivalent by every
number you currently look at; different by the number you don't'').
Fallback: report as optimizer inductive bias made measurable, without
the equivalence framing; the emblematic slot falls to the BIF alone;
strike C1(b).
\item \open Adam recast: remains an illustration unless the
$v_t$--degenerate-subspace alignment is verified; the scoped-guarantee
framing carries the honest version either way. New face from goals v5:
Adam's $\epsilon$ flagged as a candidate damping instance [VERIFY] ---
check whether the numerical-stability constant admits the
regularity-restoring reading before the flag is stated.
\item \resolved{BIF--damping relationship: the local Bayesian influence
function replaces the ill-posed Hessian inversion with SGMCMC covariance
estimation, and its leading-order term under a Laplace approximation is
exactly the classical dampened influence function, with localization
strength $\gamma$ as the dampening coefficient. Source: the Timaeus BIF
paper (Hessian-free data attribution). This is the twin arc's
load-bearing fact and the paper's cleanest
folklore-constant-becomes-principled exhibit.}
\item \partialr Susceptibilities/spectroscopy: linear-response framework
linking LLCs to susceptibilities; response-matrix separation of
functional modules (multigram, induction heads) confirmed in the
authors' expositions; pin the primary citations (Baker et al.\ 2025 and
the loss kernel / susceptibility papers) and which claims each carries.
\item \open Sampler craft (deepest assessment): temperature,
localization, weight restriction; failure modes (stuck on one branch;
mixing too fast to resolve the geometry); what \texttt{devinterp}
exposes vs.\ what remains manual. Explicitly record: a formally
stratification-preserving sampler does not exist as a citable method ---
claim the craft, flag the formal object as open (feeds \S5.4).
\item \open Worked example: rLLC on attention heads (arXiv:2410.02984)
--- reproduce or report; confirm it doubles as the \S5 interpretability
bridge.
\end{itemize}

% ============================================================================
\section{Section 5 Needs}
% ============================================================================

\begin{itemize}[itemsep=4pt]
\item \open Phase-transition computation on a DL model (may be novel; or
borrow conceptually from the applied-stats binomial mixture with an
explicit bridge).
\item \open Quasi-singular RLCT-growth claim: established, emerging, or
conjecture? Literature check before it is stated as delivered.
\item \open The specification-gap program's ancestor: pin the citation
for the optimizer that targets the free energy and demonstrably selects
lower-$\lambda$ strata (goals \S5's ``explicit ancestor already in the
literature''); confirm what it demonstrates vs.\ what the program
conjectures. The hyperparameter-dependence of final $\lambda$ within a
single algorithm is flagged in goals as an open experiment --- record it
here as a candidate follow-up, not a paper dependency.
\item \open Flag sketches, one credible symmetry argument each:
activation/representation (rotations/scalings of latent directions;
manifold-hypothesis and regular-base-flow breakage); policy/value
(behavior-preserving reparameterizations; NPG inheriting the
local-differential failure); interpretability (non-unique circuits,
superposition as degeneracy on the singular set). Bar: ``yes, that
probably extends,'' not proof. Interpretability is the developed flag
(half-page); the other two compress to a credible sentence each.
\item \open Claim 3 (the $n$-revelation) strong form: lands here as open
question if the \S2 literature check (C4) does not clear it.
\item \open \S5.4 horizon: the stratified-space sampler direction ---
one or two citations that physics handles potential functions on
singular configuration spaces (to make the ``machinery exists
elsewhere'' point credible) without claiming an MCMC literature that
does not exist.
\end{itemize}

% ============================================================================
\section{Cross-Cutting}
% ============================================================================

\begin{itemize}[itemsep=4pt]
\item \open Bibliography pass: primary sources for every audit-table row
(the table currently cites from memory plus secondary confirmation).
\item \open The companion blog extraction from \S1 (convert forward
references; end on the open question).
\item \open Tone pass against framing decision (i): the provocative
register only where the diagnosis tag is ``geometry'' AND a better
computable tool exists today (singular VI pending C2; BIF); the humble
register where the tag is ``computation.''
\item \open Notation convention with the applied-stats sibling: this
paper writes $\lambda$ replaces $p/2$; the sibling writes $d/2$. Decide
pair-wise ($p$ avoids the collision with data dimension) and apply to
whichever document yields.
\item \open Phrase-bank continuity: ``how much room the level sets
leave'' is the walk-away/\S2 phrasing that \S3's $\lambda$ definition
must echo verbatim when the volume scaling
$\mathrm{Vol}\{K \le \varepsilon\} \sim c\,\varepsilon^{\lambda}
(-\log\varepsilon)^{m-1}$ is introduced (``the room shrinks like
$\varepsilon^{\lambda}$''). A writing-consistency check, not research.
\item \resolved{Adam guarantees scoping (was: ``verify stochastic
optimization guarantees assume non-degenerate Hessian''): the accurate
statement is that nonconvex guarantees deliver $\epsilon$-stationarity
under smoothness assumptions and are silent on the singular questions;
stronger promises invoke PL/strong convexity. The old item's framing was
wrong and is retired; citation selection remains under \S1 Pillar 4
(placement now \S4.2).}
\end{itemize}

\end{document}
